% SGA 3, Expose V, Loop 1 English reconstruction -- DISJOINT FUTURE STAGING.
% Scope: component-local pp. 28--31 (combined-reader pp. 290--293),
% re-edited headers 278--281. This begins in the continuation of Lemma 7.4,
% closes Section 7, contains Section 8, and opens Section 9 through 9(b).
% Authority: Polo--Gille Exp5-13oct24.pdf, SHA-256
% 9198200633F929FE1822520371A9200DA4F8CF2513EFAB3D6A0C7E9330DB84CF.
% Jacob Reinhold's CC BY 4.0 Markdown is comparison/drafting material only.
% Integration order: only after the separately owned local pp. 24--27 tranche.

并设 \(x\) 为 \(F_2\) 的一点.

\begin{enumerate}
\item[\SGAInlineRefTarget{lem:sga3-V-7.4:item:i}(i)] 若 \(u\) 平坦, 则 \(d'\) 在 \(x\) 处平坦当且仅当
  \(d\) 在 \(v(x)\) 处平坦.
\item[(ii)] 若 \(d\) 局部有限型, 则 \(d'\) 在 \(x\) 处拟有限
  当且仅当 \(d\) 在 \(v(x)\) 处拟有限.%
  \footnote{\label{note:sga3-V-42}编者注: 由基变换可知这些条件是充分的
  (EGA~II, 6.2.4(iii), 以及
  EGA~IV\(_2\), 2.1.4). 反之, 令
  \(y=d'(x)\) 及 \(z=u(y)=d(v(x))\), 并假设 \(d'\) 在
  \(x\) 处平坦, 且 \(u\) (因而 \(v\) 也) 平坦. 则
  \(\mathcal O_{v(x)}\to\mathcal O_x\) 忠实平坦,
  \(\mathcal O_z\to\mathcal O_y\to\mathcal O_x\) 也忠实平坦. 因此,
  \(\mathcal O_z\to\mathcal O_{v(x)}\) 忠实平坦
  (EGA~IV\(_2\), 2.2.11(iv)). 最后, 假设 \(d\) 局部
  有限型, 且 \(d'\) 在 \(x\) 处拟有限. 则
  \(v(x)\) 在其纤维 \(d^{-1}(z)\) 中孤立, 因为 \(x\) 在其纤维
  \(d'^{-1}(y)=d^{-1}(z)\otimes_{\kappa(z)}\kappa(y)\) 中孤立. 因此由 Chevalley
  半连续性定理, \(v(x)\) 有一个开邻域, 其中每一点都在其纤维中孤立
  (EGA~IV\(_3\), 13.1.3 及 13.1.4), 因而 \(d\) 在
  \(v(x)\) 处拟有限.}
\end{enumerate}

由此证明了 \(X_0\) 有一个由饱和开子集 \(W\) 组成的覆盖,
使得 \(X_\ast\) 在 \(W\) 上诱导的群胚 \(W_\ast\)
具有一个拟截面.%
\footnote{\label{note:sga3-V-43}编者注: 下文经过修改,
以利用引理~\hyperref[lem:sga3-V-6.1]{6.1} 中新增的内容.}

由引理~\hyperref[lem:sga3-V-6.1]{6.1} 和定理~\hyperref[thm:sga3-V-7.1]{7.1}
之后陈述的化简, 这就证明了该定理的断言
\hyperref[thm:sga3-V-7.1-i]{(i)} 和
\hyperref[thm:sga3-V-7.1-iii]{(iii)}, 以及如下事实:
\(p\) 满射且为开映射, \(p\) 和 \(Y\to S\)
局部有限呈示. 此外, 由于 \(X_0\to S\)
拟射影, 因而分离且有限型, \(p\) 分离;
引理~\hyperref[lem:sga3-V-6.1-ii]{6.1(ii)} 的证明表明,
\(p\) 和 \(Y\to S\) 有限呈示.

为证明 \(p\) 固有, 只需再证明它普遍闭.
由于此断言在 \(Y\) 上是局部的, 可以在一个饱和开子集 \(W\) 上工作,
使得 \(X_\ast\) 在 \(W\) 上诱导的群胚
\(W_\ast\) 具有一个拟截面 \(U\)
(因为 \(X_0\) 被这样的开子集覆盖). 采用
\hyperref[par:sga3-V-6-a]{6(a)} 中的记号, 有交换图
\SGAThreeDiagram[0.48\textwidth]
  {figures/exp5/exp5_localp28_WVU_Z_commutative.png}
  {开子集 \(Z\subseteq Y\) 上的固有性图.}
  {fig:sga3-V-7.1-properness}
  {第 V 篇, 组件局部第~28 页 / 合订阅读版第~290 页 / 修订版页眉第~278 页}
其中 \(Z\) 为 \(Y\) 的开子集, 所有箭头均满射,
且 \(q\) 整. 由假设,
\(d_0:X_1\to X_0\) 固有, 故由 \(d_0\)
经基变换得到的 \(u\) 也固有. 因此 \(r\) 普遍闭,
又因 \(v\) 满射, \(p\) 也普遍闭.

最后, 由于 \(p\) 满射且普遍闭, 而
\(X_0\) 拟射影, 因而分离, 对角线
\(\Delta_{Y/S}(Y)\) 在 \(Y\times_S Y\) 中闭:
它是对角线 \(\Delta_{X_0/S}(X_0)\) 在
\(p\times p\) 下的像. 所以 \(Y\) 在 \(S\) 上分离.
这就完成了定理~\hyperref[thm:sga3-V-7.1-ii]{7.1(ii)} 的证明.

定理~\hyperref[thm:sga3-V-7.1-iv]{7.1(iv)} 中的断言
在 \(Y\) 上是局部的. 由于饱和开子集 \(U_z\) 覆盖
\(X_0\), 只需分别以 \(U_z\) 和 \(V_z=p(U_z)\)
代替 \(X\) 和 \(Y\) 后验证它们. 如定理~\hyperref[thm:sga3-V-7.1]{7.1}
证明开头所指出的, 引理~\hyperref[lem:sga3-V-1.1]{1.1},
\hyperref[sga3:exp5:lemma:1.2]{1.2} 及
\hyperref[lem:sga3-V-6.1-i]{6.1(i)} 表明,
\((V_z,p|_{U_z})\) 在 \((\mathbf{Sch})\) 中及赋环空间范畴中
都是 \(X_\ast\) 在 \(U_z\) 上诱导的群胚的余核.
\(d=(d_0,d_1)\) 为单态射这一假设
在基变换 \(U_z\to X_0\) 下保持.
因此, 定理~\hyperref[thm:sga3-V-7.1-iv]{7.1(iv)} 的前两个断言
由引理~\hyperref[lem:sga3-V-6.1-iv]{6.1(iv)} 得到.

最后证明 \hyperref[thm:sga3-V-7.1-iv]{(iv) 部分}末尾陈述的推论
(见 \(\hyperref[bib:sga3-V-Ray67a]{[\mathrm{Ray67a}]}\), 定理~1(iii)). 由假设,
群胚 \(X_\ast\) 来自一个等价关系
\(R\to X_0\times X_0\). 已经证明 \(R\) 有效
(\SGARefLink{sga3:iv:definition:3:3:2}{第~IV 篇, 3.3.2}), 且
\(p:X_0\to Y=X_0/R\) 忠实平坦且有限呈示.
因此, 若 \((M)\) 表示忠实平坦且局部有限呈示态射族,
则 \(R\) 是
\((M)\)-有效的. 由
\SGARefLink{sga3:iv:corollary:6:3:3}{第~IV 篇, 6.3.3}, 对 \((Y,p)\) 因而
表出 fppf 拓扑下 \(X_0\) 模 \(R\) 的商层,
而关于基变换的断言由
\SGARefLink{sga3:iv:proposition:3:4:3:1}{第~IV 篇, 3.4.3.1} 得到.

\section*{8. 对未必固有的平坦群胚取商}
\phantomsection\label{sec:sga3-V-8}
\addcontentsline{toc}{section}{8. 对未必固有的平坦群胚取商}

\medskip
\phantomsection\label{thm:sga3-V-8.1}
\noindent\textbf{定理 8.1.}---%
\footnote{\label{note:sga3-V-44}编者注: \(X_0\) 存在满足\hyperref[thm:sga3-V-8.1]{此定理}
结论的最大开子集 \(W\).
事实上, 设 \(W\) 为\hyperref[thm:sga3-V-8.1]{此定理}中的开子集,
\(W^\sharp\) 为 \(W\) 的稠密饱和开子集. 由于 \(p\)
为开映射, \(V^\sharp=p(W^\sharp)\) 在 \(V\) 中开,
且因 \(W^\sharp\) 饱和,
\(W^\sharp=p^{-1}(V^\sharp)\). 由
引理~\hyperref[lem:sga3-V-1.1]{1.1}, \(V^\sharp\) 是
在 \(W^\sharp\) 上诱导的群胚的余核. 因此, 满足\hyperref[thm:sga3-V-8.1]{此定理}
结论的两个开子集 \(W\) 和 \(W'\) 可以沿其交集
\(W^\sharp\) 粘合, 而条件 \hyperref[thm:sga3-V-8.1-i]{(i)}--\hyperref[thm:sga3-V-8.1-iv]{(iv)},
以及 \(p\) 和 \(V\to S\) 局部有限呈示的性质均保持.
如引理~\hyperref[lem:sga3-V-6.1-ii]{6.1(ii)} 的证明中一样,
结论 (ii') 由 \(X_0\) 在 Noether 概形
\(S\) 上有限型这一假设得到.

引理~\hyperref[lem:sga3-V-1.1]{1.1} 和
\hyperref[sga3:exp5:lemma:1.2]{1.2} 还表明, 所有满足如下条件的饱和开子集
\(U\subseteq X_0\) 的并, 是满足定理~\hyperref[thm:sga3-V-8.1]{8.1}
\hyperref[thm:sga3-V-8.1-i]{条件~(i)} 的最大饱和开子集
\(\Omega\subseteq X_0\): 开子集
\(p(U)\subseteq Y\) 是概形, 且
\(p|_U:U\to p(U)\) 是概形态射.
\hyperref[thm:sga3-V-8.1]{此定理}表明
\(\Omega\) 包含一个稠密开子集 \(W\), 但 \(\Omega\) 满足
\hyperref[thm:sga3-V-8.1-ii]{性质 (ii)}--\hyperref[thm:sga3-V-8.1-iv]{(iv)} 并非显然.
关于更精确的结果, 以及 fppf 商层
\(\widetilde{X/R}\) 的可表性 (其中 \(R\) 表示以 \(X=X_0\) 为底的群胚 \(X_\ast\)),
见 \(\hyperref[bib:sga3-V-Ray67a]{[\mathrm{Ray67a}]}\), \(\hyperref[bib:sga3-V-Ray67b]{[\mathrm{Ray67b}]}\),
以及 \(\hyperref[bib:sga3-V-An73]{[\mathrm{An73}]}\) 的附录~I.
这些文献采用更弱的假设:
\(S\) 可以任意, \(X\) 在 \(S\) 上局部有限型,
且 \(R\) 是以 \(X\) 为底的 \(S\)-群胚, 使得 \(d_0\)
(因而 \(d_1\) 也) 平坦且有限呈示. 特别地, 若
\(\widetilde{X/R}\) 由 \(S\)-概形 \(Y\) 表出, 则 \(Y\)
也是赋环空间范畴中的余核. 反之一般不成立
(见 \(\hyperref[bib:sga3-V-Mum65]{[\mathrm{Mum65}]}\), 第~0 章, \S3,
例~0.4, 引于 \(\hyperref[bib:sga3-V-Ray67a]{[\mathrm{Ray67a}]}\), 注~1),
但当 \(d=(d_0,d_1)\) 为浸入时成立. 在此假设下,
\(p:\Omega\to Z:=\Omega/R_\Omega\) 忠实平坦且有限呈示.
若此外 \(S\) 局部 Noether, 则
\(X\) 的余维一的点 \(x\) 属于 \(\Omega\) 当且仅当
在 \(\operatorname{Spec}(\mathcal O_{X,x})\) 上诱导的群胚的图闭.
见
\(\hyperref[bib:sga3-V-Ray67a]{[\mathrm{Ray67a}]}\), 命题~1;
\(\hyperref[bib:sga3-V-Ray67b]{[\mathrm{Ray67b}]}\),
命题~1 及定理~2, 1, 4; 以及 \(\hyperref[bib:sga3-V-An73]{[\mathrm{An73}]}\),
第~66--67 页的定理~5, 6 和第~49 页的命题~3.3.1.
关于代数闭域 \(k\) 上代数群的作用,
另见 \(\hyperref[bib:sga3-V-DR81]{[\mathrm{DR81}]}\).}%
\textit{设 \(S\) 为 Noether 概形, 并设}
\SGAThreeDiagram[0.57\textwidth]
  {figures/exp5/exp5_localp29_theorem81_groupoid.png}
  {定理~\hyperref[thm:sga3-V-8.1]{8.1} 中的 \((\mathbf{Sch}/S)\)-群胚.}
  {fig:sga3-V-8.1-groupoid}
  {第 V 篇, 组件局部第~29 页 / 合订阅读版第~291 页 / 修订版页眉第~279 页}
\textit{为一个 \((\mathbf{Sch}/S)\)-群胚, 使得 \(d_1\) 平坦且
有限型, \(X_0\) 在 \(S\) 上有限型, 且以 \(d_0\) 和 \(d_1\) 为分量的态射
\(X_1\to X_0\times_S X_0\)
拟有限.}

\textit{存在一个稠密饱和开子集 \(W\subseteq X_0\),
具有下列性质:}
\begin{enumerate}
\item[(i)]\phantomsection\label{thm:sga3-V-8.1-i}%
  \textit{若
  \(W_2\overset{w'_i}{\rightrightarrows}
    W_1\overset{w_j}{\rightrightarrows}W\)
  是 \(X_\ast\) 在 \(W\) 上诱导的群胚, 则
  \((w_0,w_1)\) 在
  \((\mathbf{Sch}/S)\) 中有余核 \((V,p)\); 此外,
  \((V,p)\) 是所有赋环空间的范畴中
  \((w_0,w_1)\) 的余核;}
\item[(ii)]\phantomsection\label{thm:sga3-V-8.1-ii}%
  \textit{\(p\) 满射且为开映射;}
\item[(ii')]\phantomsection\label{thm:sga3-V-8.1-iip}%
  \textit{\(p\) 和 \(V\to S\) 有限呈示;}
\item[(iii)]\phantomsection\label{thm:sga3-V-8.1-iii}%
  \textit{以 \(w_0\) 和 \(w_1\) 为分量的态射 \(W_1\to W\times_VW\)
  满射;}
\item[(iv)]\phantomsection\label{thm:sga3-V-8.1-iv}%
  \textit{若 \((d_0,d_1)\) 为等价对, 则
  \(W_1\to W\times_VW\) 是同构, 且 \(p\) 忠实
  平坦.}
\end{enumerate}

我们将证明可以选取 \(W\), 使得
\(X_\ast\) 诱导的 \((\mathbf{Sch}/S)\)-群胚 \(W_\ast\)
具有一个拟截面 (见 \hyperref[sec:sga3-V-7]{\S7}).
由引理~\hyperref[lem:sga3-V-6.1]{6.1}
即可得到定理~\hyperref[thm:sga3-V-8.1]{8.1}.

暂且假设, 对每个相对于 \(S\) 为闭点的点 \(z\in X_0\)
(见 \hyperref[sec:sga3-V-7]{\S7}), 都有一个
饱和开子集 \(W_z\), 它具有一个拟截面, 并与
\(X_0\) 中经过 \(z\) 的每个不可约分支相交. 则
补集 \(X_0-W_z\) 饱和: 其饱和化
\(d_1(d_0^{-1}(X_0-W_z))\) 开且不与 \(W_z\) 相交. 若此
补集非空, 在其中选取一个相对于 \(S\) 为闭点的点 \(z'\),
并为 \(z'\) 配备上述开子集 \(W_{z'}\).
还可以假设 \(W_{z'}\subseteq X_0-W_z\). 则 \(W_z\) 和
\(W_{z'}\) 不交, 而 \(X_\ast\) 在
\(W_z\cup W_{z'}\) 上诱导的群胚具有一个拟截面. 此过程必然终止,
因为 \(X_0\) 只有有限多个不可约分支.
因此, 只需构造 \(W_z\).

为此可以假设 \(S\) 仿射. 设 \(y\in X_1\)
满足 \(d_1(y)=z\), \(X\) 为 \(X_0\) 中
包含 \(d_0(y)\) 的仿射开子集, \(Y\) 为 \(X\) 在
\(d_0:X_1\to X_0\) 下的逆像, 并设 \(u:Y\to X\) 和 \(v:Y\to X_0\)
为 \(d_0\) 和 \(d_1\) 诱导的态射. 由于 \(X\) 仿射,
因而拟射影, 引理~\hyperref[lem:sga3-V-7.2]{7.2} 给出一个
子概形 \(F\subseteq X_0\), 使得
\(d_0^{-1}(F)\cap d_1^{-1}(z)\) 非空,
\(d_0(d_0^{-1}(F)\cap d_1^{-1}(z))\) 有限, 且
\(d_1\) 在 \(d_0^{-1}(F)\) 上的限制于 \(v^{-1}(z)\) 的各点平坦.
于是可以重新采用
引理~\hyperref[lem:sga3-V-7.3]{7.3} 的记号, 用 \(F_1\) 和 \(F_2\)
分别表示 \(F\) 在 \(X_1\) 和 \(X_2\) 中的逆像, 等等:
\SGAThreeDiagram[0.56\textwidth]
  {figures/exp5/exp5_localp30_F_groupoid.png}
  {定理~\hyperref[thm:sga3-V-8.1]{8.1} 的证明中在 \(F_2,F_1,F\) 上诱导的图.}
  {fig:sga3-V-8.1-F-groupoid}
  {第 V 篇, 组件局部第~30 页 / 合订阅读版第~292 页 / 修订版页眉第~280 页}

如 \hyperref[lem:sga3-V-7.3]{7.3} 中一样, 可以证明
\(\widetilde d_1\) 在
\(\widetilde d_1^{-1}(z)\) 的各点拟有限. 因而自然要考虑
由 \(\widetilde d_1\) 同时平坦且拟有限的点组成的
开子集 \(F'_1\subseteq F_1\). 由
引理~\hyperref[lem:sga3-V-7.4]{7.4}, \(F'_1\) 在
\(\widetilde d'_1\) 和 \(\widetilde d'_0\) 下的逆像都由
\(F_2\) 中 \(\widetilde d'_2\) 平坦且
拟有限的点组成. 因此这两个逆像相同, 且 \(F'_1\) 具有
\(\widetilde d_0^{-1}(F')\) 的形式, 其中 \(F'\subseteq F\) 为开子集
(SGA~1, VIII, \SGARefExternal{sga3:viii:corollary:4:4}{4.4}). 因此以 \(F'\) 代替 \(F\) 后,
可以假设 \(\widetilde d_1\) 拟有限且平坦. 设 \(W_z\)
为 \(\widetilde d_1(F_1)\) 的最大开子集, 使得
\(\widetilde d_1\) 在其上有限且平坦.

开子集 \(W_z\) 未必包含 \(z\), 但它包含
\(X_0\) 中经过 \(z\) 的各不可约分支的泛点.%
\footnote{\label{note:sga3-V-45}编者注: 事实上, 设 \(\eta\) 为
这样的一个泛点. 由假设,
\(\mathcal O_{X_0,\eta}\) 为 Artin 局部环, 且
\(\mathcal O_{F_1,\eta}\) 为有限生成的
\(\mathcal O_{X_0,\eta}\)-模. 因此由 SGA~1, VIII, \SGARefExternal{sga3:viii:examples:6:5}{6.5},
\(\eta\) 有一个开邻域, 使得
\(\widetilde d_1\) 在其上有限.}
由于 \(d_0\) (同样还有 \(d_1\)) 忠实平坦且有限
呈示, 因而为开映射, SGA~1, VIII, \SGARefExternal{sga3:viii:corollary:5:7}{5.7} 表明,
\(d_0^{-1}(W_z)\) 和 \(d_1^{-1}(W_z)\) 都等于
\(\widetilde d'_2(F_2)\) 的最大开子集, 使得
\(\widetilde d'_2\) 在其上有限且平坦. 如
\hyperref[lem:sga3-V-7.3]{7.3} 中一样,
\(\widetilde d_1^{-1}(W_z)\) 在 \(\widetilde d'_0\) 和
\(\widetilde d'_1\) 下的逆像因而相同. 所以
\(\widetilde d_1^{-1}(W_z)=\widetilde d_0^{-1}(U)\), 其中
\(U\subseteq F\) 为开子集, 且 \(U\) 是 \(X_\ast\)
在 \(W_z\) 上诱导的群胚的一个拟截面.

\section*{9. 消去定理 \hyperref[thm:sga3-V-7.1]{7.1} 中的 Noether 假设}
\phantomsection\label{sec:sga3-V-9}
\addcontentsline{toc}{section}{9. 消去定理 \texorpdfstring{\protect\hyperref[thm:sga3-V-7.1]{7.1}}{\SGARefLink{thm:sga3-V-7.1}{7.1}} 中的 Noether 假设}

\medskip
\phantomsection\label{par:sga3-V-9-a}%
\noindent\textbf{a)}\quad
重新采用引理~\hyperref[lem:sga3-V-6.1]{6.1}
的记号和假设, 并设
\(\pi:S'\to S\) 为任意基变换. 若 \(f:X\to Y\) 为
\(S\)-概形的态射, 用 \(f':X'\to Y'\) 表示沿
\(\pi\) 扩张底后由 \(f\) 得到的 \(S'\)-概形态射.
采用此约定, \(p':X'_0\to Y'\) 满射, 而以 \(d'_0\) 和 \(d'_1\)
为分量的态射
\[
  X'_1\longrightarrow X'_0\times_{Y'}X'_0
\]
也满射. 因此 \(Y'\) 的底集
是 \(X'_0\) 的底集模由 \(S'\)-群胚 \(X'_\ast\) 在 \(X'_0\)
上定义的等价关系所得的商.
此外, \(q':U'\to Y'\) 整且满射. 因而 \(Y'\)
的拓扑是 \(U'\) 的拓扑的商拓扑, 从而也是
\(X'_0\) 的拓扑的商拓扑; 见
\hyperref[par:sga3-V-6-c]{\S6(c)} 中的证明.

另一方面, \(U'\) 显然是
\(S'\)-群胚 \(X'_\ast\) 的一个拟截面, 故可以应用引理~\hyperref[lem:sga3-V-6.1]{6.1}.
特别地, \(X'_\ast\) 有余核 \((Y_1,p_1)\),
而 \(Y_1\) 的底拓扑空间由 \(X'_0\) 的底拓扑空间
将 \(X'_\ast\) 定义的关系下等价的点等同而得到.
因此, 典范态射
\(Y_1\to Y'\) 是同胚. 事实上, 它是普遍
同胚. 事实上, 若 \(S''\) 为 \(S'\)-概形, 设 \(Y_2\) 为
\((d_0\times_S S'',d_1\times_S S'')\) 的余核.
将上述论证应用于基变换 \(S''\to S'\) 和 \(S''\to S\), 可知
\[
  Y_2\longrightarrow Y_1\times_{S'}S''
  \qquad\text{且}\qquad
  Y_2\longrightarrow Y\times_SS''\simeq Y'\times_{S'}S''
\]
都是同胚. 因此
\[
  Y_1\times_{S'}S''\longrightarrow Y'\times_{S'}S''
\]
也是同胚.

\medskip
\phantomsection\label{par:sga3-V-9-b}%
\noindent\textbf{b)}\quad
当群胚 \(X_\ast\) \emph{局部}具有拟截面时
(见定理~\hyperref[thm:sga3-V-7.1]{7.1} 的证明),
上述论述显然同样适用.%
\footnote{\label{note:sga3-V-46}编者注: 下文作了扩充,
以突出\hyperref[thm:sga3-V-9.0]{下面的定理~9.0}.}
例如, 有\hyperref[thm:sga3-V-9.0]{下述定理}:
% SGA3-V-V2-EOL-PRESERVATION

% 翻译统计: 源行数 329; 输出行数 327 (含本统计行); 行内/行间数学片段 311; 翻译占位 0. 底本疑点另报主线程.
